\subsection{導関数の符号から増減が分かる理由}
区間内の任意の二点 $x_1<x_2$ に平均値の定理を使うと，その間のある $c$ に対して
\[
 f(x_2)-f(x_1)=f'(c)(x_2-x_1)
\]
となる。$x_2-x_1>0$ だから，区間内で $f'>0$ なら $f(x_2)>f(x_1)$ であり，$f$ は狭義に増加する。$f'<0$ なら狭義に減少する。$f'\ge0$ の場合には $f(x_2)\ge f(x_1)$ までが保証され，値が一定の部分もあり得る。

同じ議論から，次の結果も得られる。
\begin{theorem}{導関数が同じ関数の差}\label{定理:導関数同じ}
一つの区間のすべての内点で $F'=G'$ が成り立ち，端点を含む場合にはそこで両関数が連続なら，その区間で $F-G$ は定数である。特に，導関数が恒等的に $0$ なら，元の関数は定数である。
\end{theorem}

実際，$H=F-G$ とおくと $H'=0$ であり，任意の二点について $H(x_2)-H(x_1)=0$ となる。後の積分では，微分して同じ関数になる候補が，定数の違いを除いて同じであることをこの定理で確かめる。「一つの区間」という条件にも意味がある。たとえば $x<0$ で $H(x)=0$，$x>0$ で $H(x)=1$ と定めれば，定義域のすべての点で $H'=0$ だが，左右の区間に同じ定数を使うことはできない。

